%======================================================================
\section{Arquitetura do sistema}\label{sec:arquitetura}
%======================================================================
\subsection{Visão geral}

O sistema tem três tipos de nó (Figura~\ref{fig:arquitetura}). O \textbf{vestível} lê o
acelerômetro e o giroscópio, decide se houve queda e avisa o hub. O \textbf{hub do quarto}
lê os sensores fixos, roda as regras de tempo, transforma tudo em mensagens MQTT e participa
do barramento CAN da sala. A \textbf{central} recebe as mensagens, guarda o histórico,
mostra o painel e manda as notificações.

\begin{figure}[H]
\centering
\resizebox{\linewidth}{!}{% Visão geral: pulseira -> hub -> central, com os sensores do quarto e o CAN da sala.
\begin{tikzpicture}
  % Quarto
  \node[bloco=corhub, minimum width=3.4cm] (hub) at (0,0) {\textbf{Hub do quarto}\\ESP32-S3\\regras, alertas, LED};
  \node[bloco=corvest, minimum width=3.0cm] (vest) at (-5.7,0) {\textbf{Pulseira}\\ESP32-S3 + MPU-6050\\detecta a queda};
  \node[bloco=corhub, minimum width=4.6cm, font=\scriptsize] (sens) at (0,2.1)
       {portas (MC-38), presença (HC-SR501),\\temperatura (LM35 + LM358)};
  \node[bloco=corambas, minimum width=3.0cm, font=\scriptsize] (sala) at (0,-2.3) {módulos dos outros\\grupos da sala};
  \draw[sem fio=corvest] (vest) -- node[rotulo, above=1pt] {ESP-NOW} node[rotulo, below=1pt] {cifrado} (hub);
  \draw[seta=corhub] (sens) -- node[rotulo, right] {GPIO / ADC} (hub);
  \draw[thick, corambas, {Stealth[length=2mm]}-{Stealth[length=2mm]}] (hub) -- node[rotulo, right] {CAN, 250 kbit/s} (sala);
  \begin{scope}[on background layer]
    \node[draw=black!35, dashed, rounded corners, fit=(vest)(hub)(sens)(sala), inner sep=4mm,
          label={[font=\footnotesize\bfseries, text=black!60]above left:Quarto (borda)}] {};
  \end{scope}

  % Central
  \node[bloco=corcentral, minimum width=2.4cm] (mosq) at (5.4,0) {\textbf{Mosquitto}\\broker MQTT\\com ACL};
  \node[bloco=corcentral, minimum width=2.6cm] (tele) at (8.9,1.1) {Telegraf\\$\to$ InfluxDB};
  \node[bloco=corcentral, minimum width=2.6cm] (graf) at (8.9,2.6) {Grafana\\painel ``Quartos''};
  \node[bloco=corcentral, minimum width=2.6cm] (noti) at (8.9,-1.1) {notificador\\(Python)};
  \node[bloco=cinza, minimum width=2.6cm] (tg) at (8.9,-2.6) {Telegram\\do cuidador};
  \draw[seta=corhub] (hub) -- node[rotulo, above=1pt] {MQTT} node[rotulo, below=1pt] {Wi-Fi} (mosq);
  \draw[seta=corcentral] (mosq.east) -- ++(0.35,0) |- (tele.west);
  \draw[seta=corcentral] (mosq.east) -- ++(0.35,0) |- (noti.west);
  \draw[seta=corcentral] (tele) -- (graf);
  \draw[seta=cinza] (noti) -- node[rotulo, right] {internet} (tg);
  \begin{scope}[on background layer]
    \node[draw=black!35, dashed, rounded corners, fit=(mosq)(tele)(noti)(graf), inner sep=4mm,
          label={[font=\footnotesize\bfseries, text=black!60]above left:Central (Docker)}] {};
  \end{scope}
\end{tikzpicture}
}
\caption{Arquitetura geral. Azul: hub; laranja: vestível; roxo: barramento CAN da sala;
verde: central. Setas tracejadas são enlaces sem fio.}
\label{fig:arquitetura}
\end{figure}

Três princípios orientam a divisão:

\begin{enumerate}
  \item \textbf{Processamento na borda.} A queda é decidida no vestível e as regras de tempo
        rodam no hub, que publica os alertas prontos. Isso é uma diretriz da disciplina, mas
        também uma escolha de robustez: o hub acende o LED de alerta mesmo sem rede, e a
        central pode ser trocada sem mexer nas placas.
  \item \textbf{O MQTT é o contrato.} O painel, o banco, o notificador e, no futuro, um
        backend e a nuvem consomem os mesmos tópicos. O formato de cada mensagem está num
        documento único do projeto (\code{docs/protocolo.md}) e, em C, num componente
        compartilhado pelas duas placas (\code{protocol}).
  \item \textbf{Funciona sem internet.} Rede local, broker local e hora por NTP. Só o
        Telegram precisa de internet; a nuvem, quando existir, entra por um \textit{bridge}
        do Mosquitto, sem tirar o funcionamento local.
\end{enumerate}

\subsection{Dispositivos}

\begin{table}[H]
\centering
\caption{Dispositivos do sistema.}
\label{tab:dispositivos}
\small
\begin{tabular}{@{}>{\raggedright\arraybackslash}p{2.3cm}>{\raggedright\arraybackslash}p{2.6cm}>{\raggedright\arraybackslash}p{5.0cm}>{\raggedright\arraybackslash}p{4.2cm}@{}}
\toprule
Dispositivo & Hardware & Papel & Rádio \\
\midrule
Hub do quarto & ESP32-S3-DevKitC-1 na PCB v1 & sensores do quarto, regras, alertas, recepção ESP-NOW, cliente MQTT, nó CAN & Wi-Fi STA associado (sem economia de energia) + ESP-NOW + TWAI \\
Vestível & ESP32-S3-DevKitC-1 na PCB v1 & aquisição inercial, detecção de queda, envio ESP-NOW; MQTT direto e CAN como opções & STA sem associação, só para o ESP-NOW \\
Central & PC ou Raspberry Pi 4 & broker, histórico, painéis, notificação & rede local \\
\bottomrule
\end{tabular}
\end{table}

Nenhuma placa usa o modo SoftAP nem AP+STA: o ESP-NOW funciona sobre a interface STA
(Seção~\ref{sec:modos}).

\subsection{Por que estes enlaces}\label{sec:enlaces}

\paragraph{Vestível $\to$ hub: ESP-NOW.} A Tabela~\ref{tab:enlaces} compara as três opções
para a pulseira. O ESP-NOW manda um quadro em milissegundos sem associação nem conexão, o
rádio só precisa ficar ligado durante o envio, e o alarme continua funcionando dentro do
quarto mesmo com a rede Wi-Fi fora do ar. O ponto fraco é a cobertura, limitada ao raio do
hub; se os testes mostrarem pontos cegos (corredor, áreas comuns), as saídas são mais hubs
ou o Wi-Fi como reserva, que já existe como opção do vestível.

\begin{table}[H]
\centering
\caption{Opções de enlace para o vestível.}
\label{tab:enlaces}
\small
\begin{tabular}{@{}>{\raggedright\arraybackslash}p{2.6cm}>{\raggedright\arraybackslash}p{3.9cm}>{\raggedright\arraybackslash}p{3.8cm}>{\raggedright\arraybackslash}p{3.6cm}@{}}
\toprule
Critério & \textbf{ESP-NOW $\to$ hub} (escolhido) & Wi-Fi direto $\to$ broker & BLE $\to$ hub \\
\midrule
Latência do alerta & milissegundos, sem conexão & boa se associado; segundos se reconectar & depende de conexão ou anúncio \\
Bateria & rádio ligado só para enviar & precisa manter a associação & baixa, mas a pilha é pesada \\
Sem a rede Wi-Fi & funciona (alarme local) & não funciona & funciona \\
Cobertura & raio do hub & onde o Wi-Fi chega & raio do hub, menor \\
Complexidade & média: confiabilidade por nossa conta & baixa & alta; disputa o rádio do hub \\
\bottomrule
\end{tabular}
\end{table}

\paragraph{Hub $\to$ central: MQTT.} Vários consumidores ao mesmo tempo, estado retido para
quem conecta depois, \textit{Last Will} para saber que um hub caiu e QoS por criticidade.

\paragraph{CAN.} Exigência da disciplina: integra o hub aos módulos dos outros grupos na
sala. Não é caminho até a central. O vestível também tem CAN, mas só em bancada, para
testes; a pulseira de verdade não tem transceptor.

\subsection{Topologia}

A rede sem fio forma duas estrelas (Figura~\ref{fig:topologia}): em cada quarto, o vestível
fala só com o seu hub por ESP-NOW; e os hubs se associam ao ponto de acesso da ILPI como
estações Wi-Fi, chegando à central por MQTT. O barramento CAN liga os hubs aos módulos dos
outros grupos, com terminação de 120\,$\Omega$ nas duas pontas.

\begin{figure}[H]
\centering
\resizebox{0.86\linewidth}{!}{% Topologia: estrela ESP-NOW (pulseiras -> hub), estrela Wi-Fi (hubs -> AP) e barramento CAN da sala.
\begin{tikzpicture}[font=\footnotesize]
  \node[bloco=corcentral, minimum width=2.4cm] (ap) at (0,0) {ponto de acesso\\Wi-Fi da ILPI};
  \node[bloco=corcentral, minimum width=2.4cm] (cen) at (5.0,0) {central\\(PC ou Raspberry Pi)};
  \draw[thick, corcentral] (ap) -- node[rotulo, above=2pt] {Ethernet} (cen);

  \foreach \i/\x/\nome in {1/-4.6/quarto 101, 2/0/quarto 102, 3/4.6/quarto 103} {
    \node[bloco=corhub, minimum width=2.1cm] (hub\i) at (\x,-2.6) {hub\\\nome};
    \node[bloco=corvest, minimum width=1.5cm, font=\scriptsize] (v\i) at (\x-1.0,-4.3) {pulseira};
    \draw[sem fio=corvest] (v\i) -- (hub\i);
  }
  \draw[sem fio=corhub] (hub1) -- node[rotulo, sloped, above=1pt] {Wi-Fi STA (MQTT)} (ap);
  \draw[sem fio=corhub] (hub2) -- (ap);
  \draw[sem fio=corhub] (hub3) -- (ap);
  \node[font=\scriptsize, text=corvest, anchor=east] at (-5.7,-3.45) {ESP-NOW};

  % barramento CAN
  \draw[very thick, corambas] (-6.4,-5.4) -- (6.4,-5.4);
  \node[font=\scriptsize, text=corambas, anchor=west, align=left] at (6.75,-5.4) {CAN\\da sala};
  \foreach \i in {1,2,3} {\draw[thick, corambas] ($(hub\i.south)+(0.6,0)$) -- ++(0,-2.35);}
  \foreach \x/\g in {-2.3/0xD, 2.3/0xA} {
    \node[bloco=corambas, minimum width=1.6cm, font=\scriptsize] (g\g) at (\x,-6.4) {grupo \g};
    \draw[thick, corambas] (g\g.north) -- (\x,-5.4);
  }
  \foreach \x in {-6.4,6.4} {\node[draw=corambas, fill=white, inner sep=1.2pt, font=\tiny] at (\x,-5.4) {120\,$\Omega$};}
\end{tikzpicture}
}
\caption{Topologia: estrela ESP-NOW em cada quarto, estrela Wi-Fi até a central e o
barramento CAN da sala.}
\label{fig:topologia}
\end{figure}

\subsection{Como o projeto chegou aqui}\label{sec:evolucao}

O projeto começou no primeiro semestre com outra arquitetura. A Tabela~\ref{tab:evolucao}
resume as entregas e o que mudou em cada uma.

\begin{table}[H]
\centering
\caption{Linha do tempo do projeto.}
\label{tab:evolucao}
\small
\begin{tabular}{@{}l>{\raggedright\arraybackslash}p{2.4cm}>{\raggedright\arraybackslash}p{10.2cm}@{}}
\toprule
Quando & Entrega & Decisão marcante \\
\midrule
mar/26 & S1 P1 & ILPI Butantã como cenário; Modbus RTU sobre RS-485 com um Raspberry Pi mestre; Mosquitto, InfluxDB, Grafana e Telegram; queda = 2,5\,g por 100\,ms \\
mai/26 & S1 P2 & ESP32-S3-DevKitC-1; LM35 calibrado (22 pontos) e LM358 com G = 5,5; 18 ensaios do MPU-6050 \\
jul/26 & S1 P3 & CAN no lugar do RS-485; PCB de face simples com todos os sensores; firmware em C \\
ago/26 & S2 P1 & MQTT como interface aberta; CAN validado com o professor (28/08); hub e vestível \\
set/26 & S2 P2 & interface demonstrada com um simulador (18/09) \\
out/26 & S2 P2 & PCB v1 de duas camadas, única para as duas placas; revisão completa da arquitetura: regras na borda, contrato MQTT, ESP-NOW cifrado \\
nov/26 & S2 P3 Q1 & 22 métodos de detecção de queda comparados; Random Forest embarcado no lugar do limiar \\
\bottomrule
\end{tabular}
\end{table}

As mudanças de maior impacto foram: o RS-485 deu lugar ao CAN, que a disciplina passou a
pedir; as regras saíram de um Node-RED na central e foram para o hub; a pulseira deixou o
Wi-Fi e a alimentação por USB e passou a usar ESP-NOW; e o limiar fixo de queda virou só o
gatilho de um classificador.
